\documentclass[journal]{IEEEtran}
\usepackage{cite}
\usepackage{amsmath,amssymb,amsfonts}
\usepackage[unicode]{hyperref}
\hypersetup{hidelinks}
\usepackage[UTF8]{ctex}
\setlength{\parindent}{2em}

\begin{document}

\title{SMEM: A Mixture-of-Experts Routing Global-Local Framework for Spatial Multi-Omics Domain Identification}

\author{Haobin~Xu}

\maketitle

\begin{abstract}

空间多组学技术能够同时测量空间数据及多层分子信息，为解析细胞组织与组织异质性提供了新的方式，在生理病理领域至关重要。近年来，基于空间多组学的聚类分析取得了显著进展，但现有方法存在诸多不足，
大多数一阶图不足以表示复杂组织拓扑，并且会丢失长程上下文；现有方法主要依赖固定融合策略，仅限于简单模态加权，无法有效整合各组学的信息；同时，
深聚类在不平衡空间域下容易出现 cluster collapse。为此，我们提出SMEM，一种基于MoE整合的空间多组学框架。
首先，我们引入motif高阶图，丰富局部细胞拓扑；
其次，模型通过图编码器与基于Transformer的空间表征，同时捕获局部和全局特征；
最后，在信息融合阶段，引入 MoE 模块，通过动态门控，实现dynamic expert routing，将每个 spot 路由到专门化专家。
我们还设计了稳定性感知的聚类优化策略，以缓解死簇塌陷并提升训练鲁棒性。
大量实验表明，SMEM的聚类性能优于目前所有的最先进方法，并产生了更具连贯性和生物可解释性的空间域，为鲁棒的空间多组学解码与下游生物分析提供了有效方案。
\end{abstract}

\begin{IEEEkeywords}
spatial multi-omics、domain identification、motif graph、mixture of experts、deep clustering
\end{IEEEkeywords}

\section{Introduction}

\noindent\textbf{本节整体写法：}引言建议写成 8 段。写作节奏应当是：先讲空间多组学的生物学意义与任务价值，再指出现有方法的不足，随后过渡到你的方法定位，最后用 4 点贡献收束。不要在引言中堆公式，不要过早进入实现细节。

\subsection*{第 1 段：空间多组学的生物学意义}
\textbf{写什么：} [从空间组织结构、细胞异质性、微环境依赖调控切入，说明空间多组学为什么不是简单的“多模态 + 空间”，而是组织生物学和疾病研究中的关键数据范式。]

细胞的空间组织是组织功能的基本决定因素，也是理解并研究生理病理机制的关键。
单细胞RNA测序(scRNA-seq)等单组学测量技术已经发展成熟，并改变了我们对细胞类型发现与理解的方式，但其基于解离的检测破坏了原生组织上下文，致使无法有效解析组学间的空间信息；
而空间转录组学，保留了组织原位信息，但通常仍局限于单一分子层面的测量，难以同时捕获表观遗传、蛋白质和代谢等多层级调控信息，也难以完整表征细胞状态转换、细胞间通讯及微环境依赖的分子机制。
此外，不同空间转录组平台在空间分辨率、检测通量、分子覆盖度和样本兼容性之间仍需权衡，单独依赖转录层信息也难以反映转录后调控和蛋白质功能状态。~\cite{"空间多组学整合综述"}

近年来，基于空间的多组学方法开始兴起，在2022年被《自然》杂志列为七大关键技术之一,测量对象已从单一分子层扩展到同一细胞或同一组织切片上的多模态共测量，
使研究者能够同时观察细胞内分子状态与细胞间空间组织规律。
与单组学技术不同，空间多组学展现组织结构、细胞状态、分子层级与微环境依赖调控之间的关联，
通过对单细胞基因组、转录组、蛋白质组等及空间信息的联合分析，解析细胞相互作用、空间邻域和疾病相关分子回路。
目前主流的空间多组学技术大体可分为三类，第一类是测序/空间条形码型，包括基于阵列、珠子或微流控条形码的路线，代表平台有 Visium、Slide-seq、Stereo-seq、DBiT-seq，
以及在其基础上扩展出的 SPOTS、SM-Omics、Stereo-CITE-seq、MISAR-seq 等；
第二类是原位成像型，包括基于 ISS 或多重 FISH 的平台，如 Xenium、MERFISH/MERSCOPE、seqFISH、CosMx 等；
第三类则是成像质谱型。

不同类别的空间多组学技术能够为多组学分析提供全面的生物信息与空间信息。
该技术整合的计算目标为，在保留组织几何结构与局部邻域关系的前提下，学习 biologically meaningful latent space ，
并以此支持 spatial domain identification / clustering、marker discovery、cell–cell communication、pathway/GO 分析、微环境分层、疾病区域比较和调控网络推断等。
 可以看出，空间域划分是多数下游分析的前置步骤，空间聚类是测序分析的关键步骤,其核心目标是识别基因表达与组织结构具有空间关联性的区域,从而为每个捕获的细胞点或细胞精准分配空间域标签。
因此，如何正确精准地进行空间多组学整合，并进行聚类下游任务分析，就成为了目前的核心任务。

与之对应，面向组学整合的计算方法也持续发展，大致可以分为四种类别：单组学单细胞表征学习、非空间多组学整合、空间单组学整合，以及空间多组学整合。
总体而言，图模型，尤其是利用邻接关系和局部拓扑的图神经网络，在空间相关任务中已经成为最常见的建模范式之一；
而在非空间多组学整合中，则与深度生成模型、矩阵分解和概率建模等数学方法与机器学习方法共同构成主流技术路线。
在早期的单组学单细胞分析中，许多方法已经开始借助图神经网络刻画细胞间关系。
例如，scGNN 将细胞—细胞关系表示为图结构，并结合自编码器学习，用于聚类与插补；
scGCN 利用图卷积网络实现跨组织、跨平台的数据知识迁移；scGAC 则通过图注意力自编码器学习特征表示。这些方法，为后续研究，奠定了以图建模细胞关系的技术基础，
但缺点清晰。它们主要围绕组学本身信息，构建细胞图，并未显式建模组织空间结构，也并不支持多组学的整合分析，而更多是对单组学信息进行结构化表征、降噪和特征提取。

为了使用不同信息进行联合分析，研究重心进一步转向非空间多组学整合，使用联合潜空间学习、显式跨组学约束和异质特征空间对齐等方法。
***等人提出了totalVI，对 CITE-seq 的 RNA 与蛋白数据进行联合概率建模；基于此，***等人进一步提出了MultiVI，将方法推广到 RNA–ATAC或multimodal/single-modality 混等合数据中，
使缺失模态的细胞也能被纳入统一表示。在非空间多组学整合中，GLUE是中最具代表性的框架之一，引入了 knowledge-based guidance graph，显式建模不同组学层之间的 regulatory interactions。
尽管这些方法拥有跨组学对齐能力，但它们的建模对象仍是非空间细胞集合，不能充分解决 spatial domain identification。

为弥补这一缺口，一些研究[cite]开始系统引入空间坐标、组织形态和跨切片关系。与非空间组学整合不同的是，
这类方法要求,不仅要保持组织边界的连续性，还要依据空间信息分析，恢复出真实组织结构，代表方法有SpaGCN、STAGATE、GraphST、STALinger等。
这些方法提升了空间建模能力，但也存在明显的局限性，它们依赖单一组学，难以利用多组学信息（ RNA + ADT 或 RNA + ATAC ）的互补性。

随着同一组织切片上 transcriptome–proteome、transcriptome–epigenome 乃至三模态空间共测技术的发展，
如何在保留空间组织结构的同时整合异质组学信号，已成为空间多组学分析中的核心问题。
因此，空间多组学整合方法开始被探索，并取得了显著进展。
***等人提出的SpatialGlue 是 早期代表性工作之一，
它为每个模态分别构建 spatial graph 与 feature graph，
并使用双层 attention 机制，先在模态内融合信息，再在模态间聚合不同组学表征。
该网络展示了空间多组学整合相较于特征拼接，或非空间多组学整合的优势，推动了 spatially informed cross-omics integration 的发展。
后续的SpaFusion、SpaMICS等方法，均基于SpatialGlue进行了改进，其本质仍然是使用graph表示拓扑，并通过各种操作进行模态融合，进而实现空间多组学整合。

在此基础上，后续研究进一步进行了扩展。
SpatialGlue的局限在于，其仍然依赖一阶邻域结构，更偏局部图语义，长程全局上下文与高阶拓扑表达能力有限。
随后的多种方法~\cite{}在此架构上进行了改进，并取得了长足进展。
SWITCH 将问题扩展到 unpaired spatial multi-omics integration，通过 cycle-mapping 生成 pseudo-pairs；
SpatialEx/SpatialEx+ 则进一步面向 serial-section / diagonal integration，利用 histology 作为 universal anchor 跨切片整合多组学；
而最新的 SpatialCOC 强调从 spatial prior 中学习 omics-specific spatial distributions，并显式刻画模态间 nonlinear correlations。

然而，现有方法仍存在若干尚未充分解决的问题。
第一，现有方法的跨模态融合，
通常依赖 feature concatenation、shared latent factors、attention-weighted aggregation 或连续空间校正，
也缺乏面向不同 spot 的 node-wise routing 或 expert specialization 机制。
因此，在组织异质性较强、不同模态贡献随空间位置显著变化的场景中，现有融合策略难以充分刻画模态间的互补性与差异性；

第二，许多图模型仍以固定局部邻域或 pairwise neighborhood aggregation 作为主要空间先验，其表达能力受限于一阶或低阶拓扑关系，
难以同时捕捉细粒度局部结构、高阶空间依赖与长程组织上下文。
虽然已有方法开始引入多邻域、hypergraph 或 histology-derived global context，但高阶拓扑与跨组学动态融合无法很好耦合。

第三，深度空间多组学整合通常还伴随训练层面的不稳定性。
不同组学模态在维度、噪声水平、稀疏性和生物信息量上存在显著差异，容易造成强模态主导、弱模态被抑制以及小簇结构塌陷。
已有 shared/private disentanglement 或 imbalance-aware optimization 能在一定程度上缓解信息混合不清和训练贡献失衡的问题，
但尚未充分实现面向每个节点的动态专家选择与稳定聚类协同优化。因此，一个更理想的空间多组学整合框架，
应同时具备高阶空间结构建模能力、spot-specific 跨模态贡献估计能力以及 node-wise expert routing 能力，
从而在模态贡献局部变化显著、组织结构高度异质的场景下，实现更稳健、更可解释的空间组织解析。

为了实现上述目标，本文提出 SMEM，一种基于混合专家（MoE）动态路由的多模态融合与稳定深度聚类框架。
我们的方法在三个重要方面不同于现有方法。具体来说，针对第一个问题，
不同于大多数先前方法采用的固定融合方案，我们在信息融合阶段，借鉴了大语言模型（LLM）领域的混合专家架构，并且构建了四个专家模块，
进行模态感知。在该框架下，网络进行逐节点自适应的跨组学融合，不同的 spot 根据其共识与差异模式，分别由不同专家机制整合。同时，为了增强模型表达力与稳定性，我们优化了MoE的输出。
针对第二个问题，为了解决一阶邻域图的不足，SMEM使用了 motif 增强图，丰富了组学内拓扑建模，并将其与基于 Transformer 的全局上下文建模相结合。
第三，在表征融合之外，我们显式引入稳定性感知的聚类优化，以减少死簇塌陷并提升训练鲁棒性。
这些设计选择使我们的框架特别适合异质和不平衡的空间多组学数据。



\textbf{怎么写：} [先讲技术演进，再讲 biological payoff；不要一上来堆方法名或网络结构。]


\subsection*{第 2 段：任务定义——为什么 spatial domain identification / clustering 是核心任务}
\textbf{写什么：} [定义 spatial domain identification 是什么，并说明它是 marker analysis、GO/pathway 分析、cell--cell interaction 分析、microenvironment study 等后续任务的前置步骤。]

\textbf{怎么写：} [从任务重要性写，不要过早进入你自己的模型。]


\subsection*{第 3 段：为什么非空间多组学 integration 方法不够}
\textbf{写什么：} [指出 Seurat WNN、totalVI、MultiVI、MOFA+ 等方法虽然能做多组学整合，但没有显式建模空间结构，或对模态组合有限制，因此与 spatial domain identification 的问题设定存在 problem mismatch。]

\textbf{怎么写：} [这一段要写成 “problem mismatch”，不是贬低已有方法。]


\subsection*{第 4 段：SpatialGlue 的贡献与局限}
\textbf{写什么：} [先认可 SpatialGlue 的开创性：spatial graph + feature graph + within-modality attention + between-modality attention；再指出两个局限：其一，仍然主要依赖一阶邻域结构；其二，更偏局部图语义，长程全局上下文与高阶拓扑表达仍有限。]

\textbf{怎么写：} [语气必须是“先认可其开创性，再指出你的切入点”。]


\subsection*{第 5 段：SpaFusion 的贡献与局限}
\textbf{写什么：} [讲 SpaFusion 如何把 3-node motif 高阶图、graph autoencoder 和 transformer 带入 spatial multi-omics clustering；再指出它虽然解决了“局部 vs 全局”问题，但跨模态整合仍更接近静态或粗粒度融合，对模态互补、模态冲突以及不平衡聚类稳定性处理得不够显式。]

\textbf{怎么写：} [这一段是全文最关键的前置铺垫，因为你的代码本质上更接近“在 SpaFusion 主干上的增强版本”。]

\textbf{建议优先参考的上传 PDF：} [\textit{Chen 等 - 2025 - SpaFusion A multi-level fusion model for clustering spatial multi-omics data.pdf}]

\subsection*{第 6 段：SpaMICS 带来的问题意识——一致信息与互补信息}
\textbf{写什么：} [指出跨组学信号同时包含 consistent information 与 complementary information；简单拼接或加权容易混淆这两类信息。]

\textbf{怎么写：} [不要把你的方法写成 explicit subspace disentanglement；只能写成“这类工作提示我们需要更具表达力的跨组学交互建模机制”。]

\textbf{建议优先参考的上传 PDF：} [\textit{Cai 等 - 2026 - Identifying spatial domains from spatial multi-omics data using consistent and specific deep subspac.pdf}]

\subsection*{第 7 段：SpaBalance / balanced learning 带来的问题意识——为什么需要稳定优化}
\textbf{写什么：} [讲 optimization imbalance、gradient conflict、dominant modality、dead cluster 等问题，再自然引出你代码里的 cluster-usage regularization、Q-mix 和 dead-cluster reinitialization。]

\textbf{怎么写：} [不要声称你实现了 gradient conflict coordination；应当写成“本文通过 stability-aware clustering objectives 针对 optimization instability 与 cluster collapse”。]

\textbf{建议优先参考的上传 PDF：} [\textit{Cui 等 - SpaBalance Balanced Learning for Efficient Spatial Multi-Omics Decoding.pdf}；\textit{Zhu 等 - 2025 - BMCST Balanced multi-view clustering for spatially resolved transcriptomics with Mamba-driven dynam.pdf}]

\subsection*{第 8 段：你的贡献总结}
\textbf{写法要求：} [建议直接列 4 点贡献，但要写得像顶刊论文中的 contributions，不要写成项目汇报。]

\textbf{按此前建议对应的 4 点贡献框架：}
\begin{enumerate}
    \item [贡献 1：] [构建 motif 增强且经过生物学 refined 的图，以联合编码空间邻近性、表型相似性和高阶拓扑。]
    \item [贡献 2：] [设计 shared-weight dual-graph GCN 与 transformer 分支耦合的局部--全局组学内编码器，并配以自适应残差注意融合。]
    \item [贡献 3：] [引入 mixture-of-experts 跨组学融合模块，以动态路由互补与冲突的模态信号。]
    \item [贡献 4：] [开发稳定性感知的深聚类策略，以缓解 cluster collapse，并提高不平衡空间域下的鲁棒性。]
\end{enumerate}

\section{相关工作（Related Works）}

\noindent\textbf{本节整体写法：}本节不要写成“方法罗列”。建议按问题轴组织，从技术背景、双图整合、高阶图与局部--全局表示、consistent/specific 建模、balanced/robust optimization 五条线索展开，最后用一段总结 remaining gaps，把 Method 顺势引出来。

\subsection{Spatial multi-omics technologies and computational objectives}
\textbf{写什么：} [用 1 段交代 spatial multi-omics 数据的技术来源与核心计算目标：不是单纯 integration，而是在保留组织结构的前提下学习 biologically meaningful latent space。]

\textbf{怎么写：} [技术背景只写到支撑计算问题即可，不要把这一节写成实验技术综述。]

\textbf{建议优先参考的上传 PDF：} [\textit{Kiessling和Kuppe - 2024 - Spatial multi-omics novel tools to study the complexity of cardiovascular diseases.pdf}；\textit{Tang 等 - 2026 - Modal-NexT Towards unified heterogeneous cellular data integration.pdf}]

\subsection{Spatially informed dual-graph integration}
\textbf{写什么：} [重点写 SpatialGlue：spatial graph、feature graph、within-modality attention、between-modality attention。]

\textbf{怎么写：} [先总结其贡献，再指出局限：更多依赖一阶邻域，且对跨模态互补/冲突的刻画仍偏弱。]

\textbf{建议优先参考的上传 PDF：} [\textit{2024-NM-SpatialGlue.pdf}]

\subsection{High-order graph and local-global representation learning}
\textbf{写什么：} [重点写 SpaFusion：3-node motif 高阶图、graph autoencoder 提取局部结构、transformer 提取全局上下文、多层级融合与 consensus clustering。]

\textbf{怎么写：} [要明确告诉审稿人：你的方法是沿着这条路线继续发展，而不是凭空出现。]

\textbf{建议优先参考的上传 PDF：} [\textit{Chen 等 - 2025 - SpaFusion A multi-level fusion model for clustering spatial multi-omics data.pdf}]

\subsection{Consistent-specific modeling across omics}
\textbf{写什么：} [重点写 SpaMICS 的 omics-shared / omic-specific self-expression、low-rank constraint、discriminative constraint。]

\textbf{怎么写：} [这一节的作用是告诉读者：“consistent + complementary”已经是被接受的问题表述；最后一句要接回你自己的 dynamic expert routing 视角。]

\textbf{建议优先参考的上传 PDF：} [\textit{Cai 等 - 2026 - Identifying spatial domains from spatial multi-omics data using consistent and specific deep subspac.pdf}]

\subsection{Balanced and robust optimization for multi-view / multi-omics learning}
\textbf{写什么：} [写 SpaBalance 的 balanced learning、gradient conflict、shared/private learning；必要时补充 BMCST 的 multi-view imbalance / dominant-view 问题意识。]

\textbf{怎么写：} [要区分 conceptual relevance 与 direct implementation：你是借用了问题意识，而不是复现了它们的具体优化器。]

\textbf{建议优先参考的上传 PDF：} [\textit{Cui 等 - SpaBalance Balanced Learning for Efficient Spatial Multi-Omics Decoding.pdf}；\textit{Zhu 等 - 2025 - BMCST Balanced multi-view clustering for spatially resolved transcriptomics with Mamba-driven dynam.pdf}]

\subsection{Summary of remaining gaps}
\textbf{这一段建议直接概括为 4 个 gap：}
\begin{itemize}
    \item [Gap 1：] [缺少同时覆盖一阶、高阶和长程依赖的表征机制。]
    \item [Gap 2：] [缺少能够处理互补与冲突信号的动态跨模态路由机制。]
    \item [Gap 3：] [缺少面向不平衡空间域的稳定聚类策略。]
    \item [Gap 4：] [缺少把 graph construction、representation learning、cross-omics fusion 与 clustering stabilization 一体化的框架。]
\end{itemize}

\section{方法（Method）}

\noindent\textbf{整篇论文的一句话定位：} [本文应被写成一个面向复杂组织拓扑、跨模态互补/冲突信息以及不平衡空间域稳定优化的一体化框架。]

\noindent\textbf{方法定位提醒（写作时保留，成稿时删除）：} [从代码实现看，这个网络不能被写成纯复现 SpatialGlue、SpaFusion、SpaMICS 或 SpaBalance 中的任何一个；更准确的学术定位是：以 SpaFusion-style backbone 为主干，吸收 SpatialGlue 的双图融合思路，并进一步加入 MoE 跨模态路由与稳定深聚类优化。因此，正文不要把它写成 consistent-specific subspace learning 方法，也不要写成 gradient conflict coordination 方法。]

\noindent\textbf{总叙事主线（成稿时建议化为本节开头的一段总述）：}
\begin{itemize}
    \item [图构建层：] [同时使用 spatial KNN graph、feature KNN graph、基于 feature graph 的 3-node motif 高阶矩阵，以及 \texttt{spatial\_adj * feature\_adj} 的 refined spatial graph。]
    \item [组学内表示层：] [每个组学使用 refined spatial graph 上的 GCN、motif-enhanced feature graph 上的 GCN，以及原始特征矩阵上的 transformer 三分支。]
    \item [组学间融合层：] [组学间支持 variance fusion 作为稳定基线，也支持 MoE fusion 作为增强版本；MoE 中包含 FeatureGatingExpert、HadamardExpert、VarianceExpert、DifferenceExpert，并通过富输入门控实现 dynamic routing。]
    \item [聚类稳定层：] [采用三路软分配、convex mixture 的 KL guidance、dense consistency、cluster-usage regularization、center separation 和 dead-cluster reinitialization 以提升不平衡空间域下的训练稳定性。]
\end{itemize}

\subsection{Problem formulation and notations}
\textbf{写什么：} [定义两个组学 $X^{(1)} \in \mathbb{R}^{N\times d_1}$、$X^{(2)} \in \mathbb{R}^{N\times d_2}$，共享空间坐标 $L \in \mathbb{R}^{N\times 2}$，目标是学习 integrated embedding $Z \in \mathbb{R}^{N\times d}$ 并得到空间域标签。]

\textbf{怎么写：} [符号体系尽量靠近 SpaFusion / SpaMICS 的写法，避免自创过多符号。]

\textbf{建议优先参考的上传 PDF：} [\textit{Chen 等 - 2025 - SpaFusion A multi-level fusion model for clustering spatial multi-omics data.pdf}；\textit{Cai 等 - 2026 - Identifying spatial domains from spatial multi-omics data using consistent and specific deep subspac.pdf}]

\subsection{Data preprocessing and graph construction}
\textbf{写法要求：} [本小节建议拆成四个自然段或四个二级小标题，忠实贴合代码，不要额外虚构未实现的数据处理流程。]

\subsubsection{Omics-specific preprocessing}
\textbf{写什么：} [RNA：filter genes，保留 HVG 3000，normalize\_total，log1p，scale，PCA 100；ADT/Protein：CLR normalize，scale，PCA。]

\textbf{怎么写：} [这里必须忠实代码，不要额外加入未在实现中出现的 batch correction 或归一化策略。]

\subsubsection{Spatial graph and feature graph}
\textbf{写什么：} [空间图使用坐标上的 KNN；特征图使用 PCA 特征上的 KNN。]

\textbf{怎么写：} [这部分的语言可借鉴 SpatialGlue 中的 graph construction 表述。]

\subsubsection{Motif-enhanced high-order graph}
\textbf{写什么：} [强调 3-node motif 共现矩阵来自 feature graph 的显式枚举统计，并通过可学习的 $k_1, k_2$ 与 feature graph 线性组合形成高阶图。]

\textbf{怎么写：} [可借鉴 SpaFusion 的高阶图叙事，但正文表达必须如实贴合当前代码实现。]

\subsubsection{Biologically refined spatial graph}
\textbf{写什么：} [把 \texttt{spatial\_adj * feature\_adj} 明确写成 Hadamard refinement，解释其目的在于 suppress boundary noise 并保留更符合生物学语义的空间连接。]

\textbf{建议优先参考的上传 PDF：} [\textit{2024-NM-SpatialGlue.pdf}；\textit{Chen 等 - 2025 - SpaFusion A multi-level fusion model for clustering spatial multi-omics data.pdf}；\textit{Cai 等 - 2026 - Identifying spatial domains from spatial multi-omics data using consistent and specific deep subspac.pdf}]

\subsection{Local-global intra-omics representation learning}
\textbf{写什么：} [对每个组学 $v$，使用 refined spatial graph 上的 shared-weight GCN encoder、motif-enhanced feature graph 上的 shared-weight GCN encoder，以及 omics feature matrix 上的 transformer encoder，得到 $z_s^{(v)}$、$z_f^{(v)}$、$z_g^{(v)}$。]

\textbf{怎么写：} [一定要写成“shared-weight dual-graph GCN + transformer”，不要写成两个彼此独立的 GCN encoder，因为当前代码中两条 GCN 分支共享同一组编码器权重。]

\textbf{为什么这样写：} [refined spatial graph 用于保留局部组织拓扑；motif-enhanced feature graph 用于补充高阶表型相似性；transformer 用于建模长程全局依赖。]

\textbf{建议优先参考的上传 PDF：} [\textit{Chen 等 - 2025 - SpaFusion A multi-level fusion model for clustering spatial multi-omics data.pdf}；\textit{2024-NM-SpatialGlue.pdf}]

\subsection{Adaptive intra-omics embedding fusion}
\textbf{写什么：} [先定义逐节点、逐维度的可学习融合系数 $\alpha_a, \alpha_b, \alpha_c = \mathrm{Normalize}(a,b,c)$，再写 $z_i^{(v)} = \alpha_a \odot z_s^{(v)} + \alpha_b \odot z_f^{(v)} + \alpha_c \odot z_g^{(v)}$；随后写基于高阶图的局部传播；若开启全局注意，则再写 global attention residual。]

\textbf{怎么写：} [必须强调这一步不是简单加权求和，而是 node- and dimension-adaptive；同时说明 global attention 可采用 dot 或 cosine 相似度，用于显式建模长程关系。]

\textbf{建议优先参考的上传 PDF：} [\textit{2024-NM-SpatialGlue.pdf}；\textit{Chen 等 - 2025 - SpaFusion A multi-level fusion model for clustering spatial multi-omics data.pdf}]

\subsection{Mixture-of-experts cross-omics fusion}
\textbf{写什么：} [逐一写清 4 个专家：FeatureGatingExpert、HadamardExpert、VarianceExpert、DifferenceExpert；再定义门控输入为 $[z_1, z_2, z_1 \odot z_2, |z_1-z_2|]$；然后写 noisy routing 与 residual variance path。]

\textbf{怎么写：} [不要只写“引入一个 MoE 模块”，而要明确每个专家分别解决什么问题：FeatureGatingExpert 用于局部跨模态互补；HadamardExpert 用于建模同位点交互增强；VarianceExpert 提供稳定保底融合；DifferenceExpert 显式建模模态差异与冲突。]

\textbf{写作提醒：} [这一部分应被写成“dynamic cross-modal complementarity routing”，不要降格成普通 attention fusion。]

\textbf{建议优先参考的上传 PDF：} [动机层面可参考 \textit{Cui 等 - SpaBalance Balanced Learning for Efficient Spatial Multi-Omics Decoding.pdf}；实现层面请忠实代码，不要虚构不存在的跨注意力结构。]

\subsection{Reconstruction learning and two-stage optimization}
\textbf{写什么：} [写 decoder 如何重构 spatial / feature branches 的图结构，以及 transformer decoder 如何重构组学特征；再写 two-stage training：先 reconstruction pretraining，再进行 end-to-end clustering optimization。]

\textbf{怎么写：} [明确 pretraining 的作用是 stabilizing latent manifold before clustering，而不是泛泛地说“用于初始化”。]

\textbf{建议优先参考的上传 PDF：} [\textit{Chen 等 - 2025 - SpaFusion A multi-level fusion model for clustering spatial multi-omics data.pdf}]

\subsection{Stability-aware deep clustering objective}
\textbf{建议的展开顺序：}
\begin{enumerate}
    \item [步骤 1：] [KMeans initialization on pretrained fused embedding。]
    \item [步骤 2：] [定义三路 soft assignments：$Q$、$Q_1$、$Q_2$。]
    \item [步骤 3：] [由 fused branch 构造 target distribution $P$。]
    \item [步骤 4：] [解释为何采用 convex mixed KL guidance，而不是直接采用几何均值版本的 $Q$。]
    \item [步骤 5：] [写 fused embedding 与 modality-specific embeddings 的 dense consistency。]
    \item [步骤 6：] [写 cluster-usage regularization。]
    \item [步骤 7：] [若启用，再写 center separation。]
    \item [步骤 8：] [写 dead-cluster reinitialization，包括 farthest / uncertainty 两类策略。]
\end{enumerate}

\textbf{怎么写：} [这一小节的核心任务是回答“为什么你的方法比 baseline 更稳”。重点必须写清三点：为什么不用几何均值 $Q$ 而用 convex mixture $Q$；为什么需要 \texttt{cluster\_usage\_weight} 去防止小簇死亡；为什么 dead-cluster reinitialization 对不平衡空间域重要。]

\textbf{建议优先参考的上传 PDF：} [动机层面可参考 \textit{Cui 等 - SpaBalance Balanced Learning for Efficient Spatial Multi-Omics Decoding.pdf}；问题意识补充可参考 \textit{Zhu 等 - 2025 - BMCST Balanced multi-view clustering for spatially resolved transcriptomics with Mamba-driven dynam.pdf}。]

\subsection{Complexity and implementation details}
\textbf{复杂度必须诚实写出：}
\begin{itemize}
    \item [复杂度点 1：] [3-node motif 的显式枚举会带来明显计算代价。]
    \item [复杂度点 2：] [global attention 存在 $O(N^2)$ 的时间/空间压力。]
    \item [复杂度点 3：] [transformer 分支带来额外 memory overhead。]
    \item [复杂度点 4：] [当前实现更适合典型规模的 paired spatial multi-omics datasets，而不是超大规模 spot 集合。]
\end{itemize}

\textbf{写作时必须如实处理的实现提醒：}
\begin{itemize}
    \item [提醒 1：] [正式主结果不能使用 ground-truth-guided checkpoint selection。若当前仍采用 \texttt{select\_best=1} 且 \texttt{best\_metric=ari}，投稿前必须修正，否则容易被审稿人质疑使用测试标签挑选最优模型。]
    \item [提醒 2：] [\texttt{adj\_k} 命令行参数当前没有真正传入 feature graph 构建；如果后续正文要写邻居数敏感性分析，必须先修复代码，再写实验。]
    \item [提醒 3：] [写作时不要声称代码实现了 SpaMICS 的 shared/specific self-expression subspace，也不要声称实现了 SpaBalance 的 gradient conflict coordination。]
\end{itemize}

\textbf{若后续补写实验设置部分，可沿用的解释口径（本文件暂不展开实验正文）：}
\begin{itemize}
    \item [\texttt{cross\_fusion=moe}：] [说明主要性能提升来自动态跨模态专家路由，而不是普通的 variance fusion。]
    \item [\texttt{cluster\_usage\_weight=0.05} + \texttt{dead\_cluster\_reinit=1}：] [说明你已显式处理小簇死亡与不平衡空间域问题。]
    \item [\texttt{q\_mix\_alpha=0.1}：] [说明聚类 guidance 以 integrated representation 为主，仅轻度吸收单模态信息。]
    \item [\texttt{emb\_attn\_mask=0} + \texttt{emb\_attn\_sim=dot}：] [说明当前最佳配置更强调全局长程关系，而不是仅依赖邻接掩膜。]
    \item [\texttt{best\_metric=ari}：] [提醒这一点只能作为开发阶段调参线索，不能直接进入正式论文的无监督模型选择叙述。]
\end{itemize}

\vspace{0.5em}
\noindent\textbf{最后提醒：} [本文件只保留写作框架与提醒，不提供论文正文，不预填参考文献条目。正式写作时，请根据各小节给出的“建议优先参考的上传 PDF”自行整理 \texttt{.bib} 文件。]

\end{document}
